% Abstract + JEL + keywords.
\begin{abstract}
\noindent New York taxi screens suggest tips as percentages of the
fee-inclusive total, so every government surcharge on the meter raises
every suggested tip. We estimate the causal pass-through of statutory
charges into gratuities: \emph{eleven cents of each legislated
surcharge dollar becomes tip}, about four-fifths of the rate at which
riders tip the fare itself. Identification comes from six sources of
quasi-experimental variation: surcharges that switch on and off at
fixed clock times, the holiday calendar, a large repricing, a flat
fare, and the spatial boundary of the 2025 congestion toll. The
per-dollar estimate is stable across doses from \$0.50 to \$5.00 and
in both directions, and riders absorb the charges rather than
re-optimize. At the 2025 schedule, \$141 million of government charges
flow through the tip base each year, generating \$15--21 million in
gratuities that no legislature enacted; ignoring tips misstates the
incidence of these charges.
\end{abstract}

\vspace{1em}

\noindent \textbf{JEL Codes:} D91, H22, J33, R48

\vspace{0.25em}

\noindent \textbf{Keywords:} default effects, tipping, tax salience, pass-through, congestion pricing
